
\sect{Constraint Satisfaction Problems}

Let us now explore a more general class of logical inference problems and discuss probabilistic entailment within that class.
We then provide further examples based on categorical constraints.
Following Chapter~5 in \cite{russell_artificial_2021}, we now define Constraint Satisfaction Problems.

\begin{definition}\label{def:csp}
    Let there be a hypergraph $\graph=(\nodes,\edges)$ and $\extnet$ be a tensor network of boolean constraint tensors $\hypercoreofat{\edge}{\catvariableof{\edge}}$ to each $\edgein$, that is
    \begin{align*}
        \extnet = \{ \hypercoreofat{\edge}{\catvariableof{\edge}} \wcols \edge\in\edges \} \, .
    \end{align*}
    The Constraint Satisfaction Problem (CSP) to $\extnet$ is the decision whether there is a state $\catindexof{\nodes}$ such that
    \begin{align*}
        \contractionof{\extnet}{\indexedcatvariableof{\nodes}} = 1 \, .
    \end{align*}
    We say the CSP is satisfiable, when there is such a state, and unsatisfiable if not.
\end{definition}

%\subsect{Deciding Entailment on Markov Networks}
%
%Deciding entailment on Markov Networks is a general class of constraint satisfaction problems.
%Here, any factor tensor in the Markov Networks produces a constraint tensor in the respective CSP.
%
%\begin{theorem}
%    \label{the:factorReduction}
%    Let $\probof{\graph}$ be a Markov Network to the Tensor Network $\extnet=\extnetasset$ on a hypergraph $\graph=(\nodes,\edges)$. % $\secnodes\subset\nodes$ be a subset and
%%	\begin{align*}
%%		\probat{\catvariableof{\secnodes}} = \normalizationof{\{\hypercoreat{\edge} \wcols \edge\in\edges \}}{\catvariableof{\secnodes}}
%%	\end{align*}
%    For each $\edge\in\edges$ we build the factor constraint cores
%    \begin{align*}
%        \sechypercoreofat{\edge}{\catvariableof{\edge}} = \nonzerocirc\hypercoreofat{\edge}{\catvariableof{\edge}} \, .
%    \end{align*}
%    Let further $\exformulaat{\catvariableof{\secnodes}}$ be a formula depending on the variables $\secnodes$, and build $\secgraph=(\nodes,\edges\cup\{\secnodes\})$.
%    Then we have that $\probof{\graph}\models\exformula$ if and only if the constraint satisfaction problem of $\secgraph$ to the constraint tensors
%    \begin{align*}
%        \{\sechypercoreof{\edge} \wcols \edge\in\edges \} \cup \{\lnot\exformula \}
%    \end{align*}
%    is unsatisfiable.
%    %and
%    %	\[ \tilde{\probtensor}[\catvariableof{\secnodes}] = \normalizationof{\{\nonzerofunction \circ \hypercoreat{\edge} \wcols \edge\in\edges \}}{\catvariableof{\secnodes}} \]
%    %Then we have for any $\exformula$ that $\probtensor\models\exformula$ if and only if $\tilde{\probtensor}\models\exformula$.
%\end{theorem}
%\begin{proof}
%    We first show, that
%    \begin{align*}
%        \nonzerocirc\probofat{\graph}{\nodevariables} =
%        \contractionof{\{\sechypercoreof{\edge} \wcols \edge\in\edges \}}{\nodevariables} \, . %\nonzerocirc\tilde{\probtensor} \, .
%    \end{align*}
%    To this end, let $\catindexof{\nodes}\in\nodestatesof{\nodes}$ be arbitrary.
%    We have $\probofat{\graph}{\indexednodevariables}=0$ if and only if at there is an edge $\edge\in\edges$ with $\hypercoreofat{\edge}{\indexedcatvariableof{\edge}}$.
%    But this is equivalent to
%    \begin{align*}
%        \contractionof{\{\sechypercoreof{\edge} \wcols \edge\in\edges \}}{\nodevariables} \, .
%    \end{align*}
%    We thus have for any $\catindexof{\nodes}\in\nodestatesof{\nodes}$
%    \begin{align*}
%        \nonzerocirc\probofat{\graph}{\indexednodevariables} =
%        \contractionof{\{\sechypercoreof{\edge} \wcols \edge\in\edges \}}{\indexednodevariables} \, . %\nonzerocirc\tilde{\probtensor} \, .
%    \end{align*}
%
%    To continue, we have $\probof{\graph}\models\exformula$ if and only if
%    \begin{align*}
%        \contraction{\extnetat{\shortcatvariables},\lnot\exformulaat{\shortcatvariables}} = 0
%    \end{align*}
%    which is equal to
%    \begin{align*}
%        \contraction{\nonzerocirc\extnetat{\shortcatvariables},\lnot\exformulaat{\shortcatvariables}} = 0 \, .
%    \end{align*}
%    We notice that this is the unsatisfiability of the claimed Constraint Satisfaction Problem.
%
%%	We first show
%%	\begin{align}\label{eq:proofFacReduction}
%%		 \nonzerocirc\probtensor = \nonzerocirc\tilde{\probtensor} \, .
%%	\end{align}
%%	The claim follows then from \theref{the:entailmentProbToLogical}.
%%	To show \eqref{eq:proofFacReduction} let there be $\indexedcatvariableof{\secnodes}$ such that $\probtensor[\indexedcatvariableof{\secnodes}]=0$.
%%	Then for any $\indexedcatvariableof{\nodes}$ extending  $\indexedcatvariableof{\secnodes}$ we have $\contractionof{\{\hypercoreat{\edge} \wcols \edge\in\edges \}}{\indexedcatvariableof{\nodes}} = 0$ and thus also $\contractionof{\{\nonzerocirc\hypercoreat{\edge} \wcols \edge\in\edges \}}{\indexedcatvariableof{\nodes}} = 0$ and $\tilde{\probtensor}[\indexedcatvariableof{\secnodes}]=0$.
%%	One can similarly show, that when $\tilde{\probtensor}[\indexedcatvariableof{\secnodes}]=0$ then also ${\probtensor}[\indexedcatvariableof{\secnodes}]=0$.
%%	The support of the distributions $\probtensor$ and $\tilde{\probtensor}$ is thus identical and \eqref{eq:proofFacReduction} holds.
%\end{proof}
%
%% Consequence: Reduction of probabilitic entailment to logical entailment.
%For any positive tensor $\hypercore$ we have
%\[ \nonzerocirc\hypercoreat{\catvariableof{\edge}} = \onesat{\catvariableof{\edge}} \, , \]
%which does not influence the distribution and can be omitted from the Markov Network.
%By \theref{the:factorReduction}, when deciding entailment, we can reduce all tensors of a Markov Network to their support and omit those with full support.
%Since the support indicating tensors $\nonzerocirc\hypercoreat{\catvariableof{\edge}}$ are boolean, each is a propositional formula and the Markov Network is turned into a Knowledge Base of their conjunctions.
%Deciding probabilistic entailment is thus traced back to logical entailment.
%
%%Exponential families
%Exponential families have a tensor network representation by a Markov Network (see \theref{the:expFamilyTensorRep}).
%However, all factors corresponding with a coordinate of the statistic $\sstat$ have a trivial support, and therefore do not influence the support of the distribution.
%The only tensors with non-trivial support are those to the boolean base measure $\basemeasure$.


\subsect{Categorical Constraints}\label{sec:categoricalTN}

%% Categorical variables with more possibilities
We so far in this chapter made the assumption that all categorical variables in factored systems to be represented by propositional logics take binary values (i.e. $\catdim=2$).
In cases where a categorical variable $\catvariable$ takes multiple values we define for each $\catindex$ an atomic formula $\catvariableof{\catindex}$ representing whether $\catvariable$ is assigned by $\catindex$ in a specific state.
%\[ \catvariableof{\catindex} =  (\catvariable = \catindex \, . \] Confusing notation
Following this construction we have the constraint that exactly one of the atoms $\catvariableof{\catindex}$ is $1$ at each state.

%% Capture constraint
%To capture the constraints resulting from this construction we introduce auxiliary parts. % of Bayesian Propositional Networks.
%Such constraints can also be expressed by a formula but would result in an unnecessary large tensor network.


%% Categorical selection map
\begin{definition}[Categorical Constraint and Atomization Variables]\label{def:catConAtomVar}
    Given a list $\catvariableof{0},\ldots,\catvariableof{\catdim-1}$ of boolean variables and a categorical variable $\catvariable$ with dimension $\catdim$ a categorical constraint is a tensor $\categoricalcoreat{\catvariableof{[\catdim]},\catvariable}$ with coordinates
    \begin{align*}
        \categoricalcoreat{\indexedcatvariableof{[\catdim]},\indexedcatvariable}
%		 \categoricalmap(\catindex,\catindexof{\variableset})
        = \begin{cases}
              1 & \ifspace \shortcatindices = \onehotmapof{\catindex} \quad \Big(\text{i.e. }\forall \catenumerator\in[\catdim]: \big(\catindex=\catenumerator\big) \Leftrightarrow \big(\catindexof{\catenumerator}=1\big) \Big)\\
              0 & \text{else} \, .
        \end{cases}
    \end{align*}
    We then call the variables  $\catvariableof{0},\ldots,\catvariableof{\catdim-1}$ the atomization variables to the categorical variable $\catvariable$.
\end{definition}

%% Decomposition
%We notice that the categorical constraint tensor is the basis encoding of the map $\categoricalmap \defcols [\atomorder] \rightarrow \bigtimes_{\atomenumeratorin} [2]$ defined as
%\begin{align*}
%    \categoricalmapat{\atomenumerator} = \onehotmapof{\atomenumerator} \, .
%\end{align*}
We notice that the categorical constraint tensor is the basis encoding of the one-hot map on $[\catdim]$.
With \theref{the:functionDecompositionBasisCP} the basis encoding $\categoricalcore$ decomposes in a basis $\cpformat$ format (see \figref{fig:CategoricalDecomposition}b) of if its coordinate maps $\categoricalmapof{\catenumerator}$, where $\catenumerator\in[\catdim]$, defined as
\begin{align*}
    \categoricalmapofat{\atomenumerator}{\secatomenumerator}
    = \begin{cases}
          1 & \ifspace \seccatenumerator=\atomenumerator \\
          0 & \text{else} \, .
      \end{cases}
\end{align*}
%
%\begin{align*}
%    \categoricalmapofat{\catenumerator}{\indexedcatvariableof{\catenumerator},\indexedcatvariable}
%    = \begin{cases}
%          1 & \ifspace \catindex=\catenumerator \andspace \catindexof{\catenumerator} = 1 \\
%          1 & \ifspace \catindex\neq\catenumerator \andspace \catindexof{\catenumerator} = 0 \\
%          0 & \text{else} \, .
%    \end{cases}
%\end{align*}
Their basis encoding are decomposed as
\begin{align}
    \categoricalcoreofat{\catenumerator}{\catvariableof{\catenumerator},\catvariable}
    = \tbasisat{\catvariableof{\catenumerator}} \otimes \onehotmapofat{\catenumerator}{\catvariable}
    + \onehotmapofat{0}{\catvariableof{\catenumerator}} \otimes \big(\onesat{\catvariable}- \onehotmapofat{\catenumerator}{\catvariable}\big) \, .
\end{align}
We thus have by \theref{the:functionDecompositionBasisCP}
\begin{align*}
    \bencodingofat{\categoricalmap}{\catvariableof{[\catdim]},\catvariable}
    = \contractionof{\left\{\categoricalcoreofat{\catenumerator}{\catvariableof{\catenumerator},\catvariable} \wcols \catindex\in[\catdim]\right\}}{\catvariable, \catvariableof{0}, \ldots, \catvariableof{\catdim-1}} \, .
\end{align*}


In the next theorem we show how a categorical constraint can be enforced in a tensor network by adding the tensor $\categoricalmap$ to a contraction.

\begin{theorem}
    For any tensor $\hypercoreat{\shortcatvariables}$ and a categorical constraint defined by an ordered subset $\catvariableof{\variableset}\subset\shortcatvariables$, a variable $\catvariable\in\shortcatvariables$ we have
    \begin{align*}
        &\contractionof{\hypercoreat{\shortcatvariables},\categoricalcoreat{\catvariableof{\variableset},\catvariable}}{\indexedcatvariables} \\
        & \quad  = \begin{cases}
              \hypercoreat{\indexedcatvariables} & \ifspace \exists \catindex : \catindexof{\variableset} = \onehotmapof{\catindex} \\
              0 & \text{else} \, .
        \end{cases}
    \end{align*}
    Here by $\catindexof{\variableset}$ we denote the restriction of $\shortcatindices$ on the set $\variableset$.
\end{theorem}
\begin{proof}
    For any $\catindexof{[\atomorder]}$ we have
    \[ \contractionof{\{\hypercoreat{\shortcatvariables},\categoricalcoreat{\catvariableof{\variableset},\catvariable}\}}{\indexedcatvariables}  =
    \hypercoreat{\shortcatindices} \cdot \categoricalcoreat{\indexedcatvariableof{\variableset},\indexedcatvariable} \, .
    \]
    If $\catindexof{\variableset} = \onehotmapof{\catindex}$ we have $\categoricalmap[\indexedcatvariableof{\variableset},\indexedcatvariable] = 1$ and thus
    \[ \contractionof{\hypercoreat{\shortcatvariables},\categoricalcoreat{\catvariableof{\variableset},\catvariable}}{\indexedcatvariables}  =  \hypercoreat{\shortcatindices}  \, . \]
    If this is not the case then $\categoricalcoreat{\indexedcatvariableof{\variableset},\indexedcatvariable} = 0$ and
    \begin{align*}
        \contractionof{\hypercoreat{\shortcatvariables},\categoricalcoreat{\catvariableof{\variableset},\catvariable}}{\indexedcatvariables}  = 0 \, . & \qedhere
    \end{align*}
\end{proof}

\begin{figure}[t]
    \begin{center}
        \begin{tikzpicture}[scale=0.35, thick] % , baseline = -3.5pt

    \begin{scope}[shift={(-15,2)}]

        \node[anchor=center] (text) at (-1,3) {${a)}$};

        \drawvariablenode{0}{0}{$\catvariableof{0}$}{T1}
        \drawvariablenode{3}{0}{$\catvariableof{1}$}{T2}
        \drawvariablenode{9}{0}{$\catvariableof{\atomorder-1}$}{T3}
        \drawvariablenode{4.5}{-5}{$\catvariable$}{C}

        \drawtextnode{6}{0}{\corelabelsize $\cdots$}

        \draw[->-] (C) -- (T1);
        \draw[->-] (C) -- (T2);
        \draw[->-] (C) -- (T3);

    \end{scope}

    \node[anchor=center] (text) at (-1,5) {${b)}$};


    \drawatomindices{0}{2}
    \draw (-1,1) rectangle (5,-1);
    \node[anchor=center] (text) at (2,0) {\corelabelsize $\categoricalcore$};
    \draw[->-] (2,-3) -- (2,-1) node[midway,left] {\colorlabelsize $\catvariable$};

    \node[anchor=center] (text) at (7,0) {${=}$};


    \begin{scope}[shift={(10,2)}]

        \newcommand{\conposseldec}{4.5,-5.5}

        \draw[fill] (\conposseldec) circle (\dotsize);
        \draw (\conposseldec) -- (4.5,-7.5) node[midway, right] {\colorlabelsize $\catvariable$};

        \draw[-<-] (0,1) -- (0,-1) node[midway,left] {\colorlabelsize $\catvariableof{0}$};
        \draw (-1,-1) rectangle (1, -3);
        \node[anchor=center] (text) at (0,-2) {\corelabelsize $\categoricalcoreof{0}$};
        \draw[-<-] (0,-3) to[bend right=20] (\conposseldec);


        \draw[-<-] (3,1) -- (3,-1) node[midway,left] {\colorlabelsize $\catvariableof{1}$};
        \draw (2,-1) rectangle (4, -3);
        \node[anchor=center] (text) at (3,-2) {\corelabelsize $\categoricalcoreof{1}$};
        \draw[-<-] (3,-3) to[bend right=20]  (\conposseldec);

        \node[anchor=center] (text) at (6,-2) {$\cdots$};

        \draw[-<-] (9,1) -- (9,-1) node[midway,left] {\colorlabelsize $\catvariableof{\atomorder-1}$};
        \draw (7.75,-1) rectangle (10.25, -3);
        \node[anchor=center] (text) at (9,-2) {\corelabelsize $\categoricalcoreof{\atomorder-1}$};
        \draw[-<-] (9,-3) to[bend left=20]  (\conposseldec);


    \end{scope}


\end{tikzpicture}
    \end{center}
    \caption{Representation of a categorical constraint in a $\cpformat$ Format tensor network.
    a) Representation of the dependency of the graphical model.
    b) Tensor Representation with further network decomposition.
    }
    \label{fig:CategoricalDecomposition}
\end{figure}

% Drop?
\begin{remark}[Constraint Satisfaction Problems of Categorical Constraints]
    We can define CSPs by collection of categorical constraints.
    An example, where the corresponding Constraint Satisfaction Problem is unsatisfiable are the categorical constraints to the three sets
    \[ \{\catvariableof{0},\catvariableof{1},\catvariableof{2},\catvariableof{3}\} \, , \, \{\catvariableof{0},\catvariableof{1}\}\, ,\,\{\catvariableof{2},\catvariableof{3}\} \, . \]
%	Besides the categorical cores also the datacores have a similar bayesian network affecting the atoms by another hidden variable.
%	Combining both is well-defined, only when all datapoints satisfy the categorical constraints (that is only one of the atoms in each constraint is active).
\end{remark}


\begin{example}[Sudoku]
    \label{exa:sudoku}
    An interesting example, where categorical constraints are combined is Sudoku, the game of assigning numbers to a grid (see for example Section~5.2.6 in \cite{russell_artificial_2021}).
    The basic variables therein are $\catvariableof{i,j}$, with $\catdimof{i,j}=n^2$ and $i,j\in[n^2]$.
    By understanding $i$ as a line index and $j$ as a column index, they are ordered in a grid as sketched in \figref{fig:sudokuGrid} in the case $n=3$.

    For a $n\in\nn$ we further define the atomization variables $\catvariableof{i,j,k}$ where $i,j,k\in[n^2]$ and $\catdimof{i,j,k}=2$.
    These $n^6$ variables are the booleans indicating whether a specific position has a specific number assigned.
    The consistency of the atomization variables to the basic variables is then for each $i,j\in[n^2]$ ensured by the categorical constraints on the sets
    \[ \{\catvariableof{i,j,k} \wcols k\in [n^2] \} \, . \]

    We further have $3\cdot n^2$ constraints by the
    \begin{itemize}
        \item Row constraints: Each number $k$ appears exactly once in each row $i\in[n^2]$, captured by the constraints
        \[ \{\catvariableof{i,j,k}  \wcols j \in [n^2] \} \, . \]
        \item Column constraints: Each number $k$ appears exactly once in each column $j\in[n^2]$, captured by the constraints
        \[ \{\catvariableof{i,j,k}  \wcols i \in [n^2] \} \, . \]
        \item Square constraints: Each number appears exactly once in each square $s,r\in[n]$, captured by the constraints
        \[ \{\catvariableof{i+n\cdot s,j+n\cdot r,k}  \wcols i,j \in [n] \} \, . \]
    \end{itemize}

    In total we have $3\cdot n^2 + n^4$ constraints for $n^6$ variables.

    Deciding whether a Sudoku has a solution is a Constraint Satisfaction Problem \cite{simonis_sudoku_2005}, which is NP-hard \cite{agerbeck_multi-agent_2008}.
    Let us notice, that due to this large number of variables and constraints, direct solution of the problem by a global contraction is not feasible.
    For efficient algorithmic solutions, we instead refer to \secref{subsec:LocalEntailment}.

%    \begin{figure}
%        \begin{center}
%            \input{PartII/tikz_pics/network_representation/sudoku_grid.tex}
%        \end{center}
%        \caption{
%            Sudoku grid of basic categorical variables $\catvariableof{i,j}$, here drawn in the standard case of $n=3$, each with dimension $\catdim=n^2=9$.
%            Each basic categorical variables has $n^2$ corresponding atomization variables, which are further atomization variables to the row, column and squares constraints.
%            Instead of depicting those constraints by hyperedges in a variable dependency graph, we here just indicate their existence through row, column and squares blocks.
%        }\label{fig:sudokuGrid}
%    \end{figure}
\end{example}